% This is part of Mes notes de mathématique
% Copyright (c) 2008-2020
%   Laurent Claessens
% See the file fdl-1.3.txt for copying conditions.

%+++++++++++++++++++++++++++++++++++++++++++++++++++++++++++++++++++++++++++++++++++++++++++++++++++++++++++++++++++++++++++ 
\section{Changement de base}
%+++++++++++++++++++++++++++++++++++++++++++++++++++++++++++++++++++++++++++++++++++++++++++++++++++++++++++++++++++++++++++

Soit un espace vectoriel \( V\) muni de deux bases \( (e_i)_{i=1,\ldots, n}\) et \( (f_{\alpha})_{\alpha=1,\ldots, n}\). Les deux bases sont liées entre elles par
\begin{equation}        \label{EQooFRQRooSMsQQB}
    f_{\alpha}=\sum_iQ_{i\alpha}e_i.
\end{equation}
Ici \( Q\) n'est pas une application linéaire \( V\to V\) : \( Q\) est seulement un tableau de nombres, donnant les coordonnées des vecteurs \( f_{\alpha}\) dans la base de \( e_i\). Éventuellement \( Q\) peut être vu comme une application linéaire \( \eK^n\to \eK^n\).

Dans la suite nous nommerons \( Q^{-1}\) la matrice inverse de \( Q\). Inverse au sens des bêtes tableaux de nombres, sans interprétations en tant qu'application linéaire. De même pour \( Q^t\) qui est la transposée de \( Q\).

%---------------------------------------------------------------------------------------------------------------------------
\subsection{Changement de base : vecteurs de base}
%---------------------------------------------------------------------------------------------------------------------------

\begin{lemma}       \label{LEMooIHZGooOZoYZd}
    Soit un espace vectoriel \( V\) sur \( \eK\) ainsi que deux bases \( (e_i)_{i=1,\ldots, n}\), \( (f_{\alpha})_{\alpha=1,\ldots, n}\) de \( V\) liées par \( f_{\alpha}=\sum_iQ_{i\alpha}e_i\). Alors
    \begin{equation}    \label{EQooZQPAooAbKAdg}
        e_i=\sum_{\alpha}Q^{-1}_{\alpha i}f_{\alpha}.
    \end{equation}
\end{lemma}

\begin{proof}
    Nous multiplions l'égalité \( f_{\alpha}=\sum_iQ_{i\alpha}e_i\) par le nombre\footnote{Attention à la bonne interprétation de ce nombre : on fait bien référence à l'élément situé en \( (\alpha, j) \) de la matrice \( Q^{-1} \), et pas autre chose.} \( Q^{-1}_{\alpha j}\in \eK\)  et nous sommons sur \( \alpha\) :
    \begin{equation}
        \sum_{\alpha}Q^{-1}_{\alpha j}f_{\alpha}=\sum_{i\alpha}(A_{i\alpha}Q^{-1}_{\alpha j})e_i=e_j.
    \end{equation}
\end{proof}

%---------------------------------------------------------------------------------------------------------------------------
\subsection{Changement de base : coordonnées}
%---------------------------------------------------------------------------------------------------------------------------

\begin{proposition}     \label{PROPooNYYOooHqHryX}
    Soit un espace vectoriel \( V\) sur \( \eK\). Soient deux bases \( (e_i)_{i=1,\ldots, n}\) et \( (f_{\alpha})_{\alpha=1,\ldots, n}\) liées par \( f_{\alpha}=\sum_iQ_{i\alpha}e_i\). Nous considérons un même vecteur dans les deux bases : \( \sum_ix_ie_i=\sum_{\alpha}y_{\alpha}f_{\alpha}\). Alors
    \begin{enumerate}
        \item       \label{ITEMooIBAEooNaUnPD}
            \( y_{\alpha}=\sum_iQ^{-1}_{\alpha i}x_i\)
        \item       \label{ITEMooKPWTooMwdbPu}
            $x_i=\sum_{\alpha}Q_{i\alpha}y_{\alpha}$.
    \end{enumerate}
\end{proposition}

\begin{proof}
    Soit un vecteur \( x\in V\). Il peut être écrit dans les deux bases :
    \begin{equation}        \label{EQooGSJMooGQstMx}
        x=\sum_ix_ie_i=\sum_{\alpha}y_{\alpha}f_{\alpha}.
    \end{equation}
    En remplaçant \( e_i\) par sa valeur \eqref{EQooZQPAooAbKAdg} nous avons l'égalité
    \begin{equation}
        \sum_{i\alpha}x_iQ^{-1}_{\alpha i}f_{\alpha}=\sum_{\alpha}y_{\alpha}f_{\alpha}.
    \end{equation}
    Vu que les \( f_{\alpha}\) sont linéairement indépendants, l'égalité des sommes donne l'égalité de chacun de termes :
    \begin{equation}
        y_{\alpha}=\sum_ix_iQ^{-1}_{\alpha i}.
    \end{equation}
    En identifiant \( x\in V\) au vecteur dans \( \eK^n\) de ses coordonnées dans la base \( \{ e_i \}\) nous pouvons écrire
    \begin{equation}
        y_{\alpha}=(Q^{-1}x)_{\alpha},
    \end{equation}
    Le point \ref{ITEMooIBAEooNaUnPD} est prouvé.

    En ce qui concerne le point \ref{ITEMooKPWTooMwdbPu}, nous repartons encore de \eqref{EQooGSJMooGQstMx}, mais nous y substituons la définition des \( f_{\alpha}\) :
    \begin{equation}
        \sum_{i}x_ie_i=\sum_{\alpha i}y_{\alpha}Q_{i\alpha}e_i.
    \end{equation}
    Vous voulez des détails ? Allez, une étape de plus que le strict nécessaire : nous écrivons
    \begin{equation}
        \sum_i\big( x_i-\sum_{\alpha}y_{\alpha}Q_{i\alpha} \big)e_i=0.
    \end{equation}
    Par linéaire indépendance des \( e_i\), nous avons annulation de tous les coefficients, c'est à dire
    \begin{equation}
        x_i=\sum_{\alpha}Q_{i\alpha}y_{\alpha},
    \end{equation}
    comme annoncé.
\end{proof}

\begin{normaltext}
    Attention à l'ordre des indices dans la dernière égalité : la matrice \( Q\) vient avec les indices dans l'ordre \( i\alpha\), tandis que la matrice \( Q^{-1}\) vient avec les indices dans l'ordre opposé : \( \alpha i\). C'est pour cela qu'il est intéressant de noter avec des lettres latines les indices se rapportant à la première base et avec des lettres grecques ceux se rapportant à la seconde base.
\end{normaltext}

\begin{normaltext}      \label{NORMooNWKZooPMwYTO}
    Les formules de changement de coordonnées de la proposition \ref{PROPooNYYOooHqHryX} s'écrivent souvent de la façon suivante :
    \begin{enumerate}
            \item       \label{ITEMooLHQCooBRvSlp}
                \( y_{\alpha}=(Q^{-1}x)_{\alpha}\)
            \item       \label{ITEMooNXUGooJIeoBf}
                \( y=Q^{-1}x\).
            \item       \label{ITEMooEFILooNENamW}
            $x_i=(Qy)_i $
            \item       \label{ITEMooMOKHooFEJvIW}
            $x=Qy$
    \end{enumerate}
    Ces égalités reposent sur un petit paquet d'abus de notations qu'il convient de bien comprendre. Ici, \( x\) et \( y\) sont les éléments de \( \eK^n\) donnés par les composantes de \( x\) dans les bases \( \{ e_i \}\) et \( \{ f_{\alpha} \}\), et \( Q\) est vu comme une matrice, un opérateur linéaire sur \( \eK^n\). Autrement dit, le choix des bases permet d'identifier \( V\) avec \( \eK^n\) et la matrice \( Q\) avec l'application linéaire \( f_Q\) de la proposition \ref{PROPooGXDBooHfKRrv}.
\end{normaltext}

%---------------------------------------------------------------------------------------------------------------------------
\subsection{Changement de base : matrice d'une application linéaire}
%---------------------------------------------------------------------------------------------------------------------------

\begin{proposition}     \label{PROPooNZBEooWyCXTw}
    Soit une application linéaire \( t\colon V\to V\) de matrices \( A\) et \( B\) dans les bases \( \{ e_i \}\) et \( \{ f_{\alpha} \}\). Si les bases sont liées par
    \begin{equation}
        f_{\alpha}=\sum_iQ_{i\alpha}e_i,
    \end{equation}
    alors les matrices \( A\) et \( B\) sont liées par
    \begin{equation}
        B=Q^{-1}AQ.
    \end{equation}
\end{proposition}

\begin{proof}
    L'hypothèse sur le fait que \( A\) et \( B\) sont les matrices de \( t\) signifie que pour tout \( x\in V\),
    \begin{equation}
        t(x)=\sum_{ij}A_{ji}x_ie_j=\sum_{\alpha\beta}B_{\alpha\beta}y_{\beta}f_{\alpha}.
    \end{equation}
    En remplaçant \( e_j\) par son expression \eqref{EQooZQPAooAbKAdg} en termes des \( f_{\alpha}\) et \( x_i\) par son expression \( x_i=(Qy)_i\) (proposition \ref{PROPooNYYOooHqHryX}), nous avons
    \begin{subequations}
        \begin{align}
            (By)_{\alpha}&=\sum_{ij\alpha}A_{ji}(Qy)_iQ^{-1}_{\alpha j}f_{\alpha}\\
            &=\sum_{i \alpha}(Q^{-1}A)_{\alpha i}(Qy)_if_{\alpha}\\
            &=\sum_{\alpha}(Q^{-1} AQy)_{\alpha}f_{\alpha}.
        \end{align}
    \end{subequations}
    Vu que les \( f_{\alpha}\) forment une base nous en déduisons \( Q^{-1}AQy=By\). Et vu que \( y\) est un élément quelconque de \( \eK^n\), nous en déduisons l'égalité de matrices
    \begin{equation}        \label{ooWKTYooOJfclT}
        B=Q^{-1}AQ.
    \end{equation}
\end{proof}
Il s'agit bien d'une égalité de matrices, ou à la limite d'applications linéaires sur \( \eK^n\), et non d'une égalité d'application linéaire sur \( V\).

%++++++++++++++++++++++++++++++++++++++++++++++++++++++++++++++++++++++++++++++++++++++++++++++++++++++++++++++++++++++++++++++++++++++++
\section{Espaces de polynômes}
%++++++++++++++++++++++++++++++++++++++++++++++++++++++++++++++++++++++++++++++++++++++++++++++++++++++++++++++++++++++++++++++++++++++++
\label{SecEspacePolynomes}

Attention : les polynômes en soi sont définis par la définition~\ref{DEFooFYZRooMikwEL}.

Pour chaque $k>0$ donné nous définissons
\begin{equation}
\mathcal{P}_\eR^k=\{p:\eR\to \eR\,|\, p : x\mapsto a_0+a_1 x +a_2 x^2 + \cdots+a_k x^k, \, a_i\in\eR,\,\forall i=0,\ldots,k\}.
\end{equation}
Il est facile de se convaincre que la somme de deux polynômes de degré inférieur ou égal à $k$ est encore un polynôme de degré inférieur ou égal à $k$. En outre il est clair que la multiplication par un scalaire ne peut pas augmenter le degré d'un polynôme. L'ensemble $\mathcal{P}_\eR^k$ est donc un espace vectoriel muni des opérations héritées de $\mathcal{P}_{\eR}$.

La base canonique de l'espace $\mathcal{P}_\eR^k$ est donné par les monômes $\mathcal{B}=\{x\mapsto x^j \,|\, j=0, \ldots, k\}$. Le fait que cela soit une base est vraiment facile à démontrer et est un exercice très utile si vous ne l'avez pas encore vu dans un cours précédent.

Nous allons maintenant étudier trois applications linéaires de $\mathcal{P}_\eR^k$ vers des autres espaces vectoriels
\begin{description}
  \item[L'isomorphisme canonique  $\phi:\mathcal{P}_\eR^k \to\eR^{k+1}$] Nous définissons $\phi$ par les relations suivantes
\[
\phi(x^j)=e_{j+1}, \qquad \forall j\in\{0,\dots, k\}.
\]
Cela veut dire que pour tout $p$ dans $\mathcal{P}_\eR^k$, avec $p(x)=a_0+a_1 x +a_2 x^2 + \cdots+a_k x^k$, l'image de $p$ par $\phi$ est
\[
\phi(p)=\phi\left(\sum_{j=0}^k a_j x^j\right)=\sum_{j=0}^k a_j e_{j+1}.
\]
\begin{example} Soit $k=5$ on a
  \begin{equation}
    \phi(-8-7x-4x^2+4x^3+2x^5)=
  \begin{pmatrix}
    -8\\
    -7\\
    -4\\
    4\\
    0\\
    2
  \end{pmatrix}.
  \end{equation}
\end{example}

Cette application est clairement bijective et respecte les opérations d'espace vectoriel, donc elle est un isomorphisme d'espaces vectoriels. L'existence d'un isomorphisme entre $\mathcal{P}_\eR^k$  et $\eR^{k+1}$ est un cas particulier du théorème qui dit que  pour chaque $m$ dans $\eN_0$ fixée, tous les espaces vectoriels sur $\eR$ de dimension $m$ sont isomorphes à $\eR^m$. Vous connaissez peut être déjà ce théorème depuis votre cours d'algèbre linéaire.
    \item[La dérivation $d: \mathcal{P}_\eR^k \to \mathcal{P}_\eR^{k-1}$] L'application de dérivation $d$ fait exactement ce qu'on s'attend d'elle
\[
d(x^0)=d(1)=0, \qquad d(x^j)=j x^{j-1}, \quad \forall j\in\{1,\dots, k\}.
\]
Cette application n'est pas injective, parce que l'image de $p$ ne dépend pas de la valeur de $a_0$, donc si deux polynômes sont les mêmes à une constante près ils auront la même image par $d$.

\begin{example} Soit $k=3$ on a
  \begin{equation}
    d(-8-12x+4x^3)= -12 (1) + 4 (3x^2) = -12+12 x^2.
    \end{equation}

    Noter que $d(-30-12x+4x^3)=d(-8-12x+4x^3)$. Cela confirme, comme mentionné plus haut, que la dérivée n'est pas injective.
\end{example}
      \item[L'intégration $I: \mathcal{P}_\eR^k \to \mathcal{P}_\eR^{k+1}$] Nous pouvons définir une application que est <<à une constante prés>> l'application inverse de la dérivation. Cette application est définie sur les éléments de base par
          \begin{equation}
                I(x^j)= \frac{x^{j+1}}{j+1}.
          \end{equation}
          Bien entendu la raison d'être et la motivation de cette définition apparaîtra lorsque nous développerons une théorie générale de l'intégration.

\begin{example}
   Soit $k=4$ on a
  \begin{equation}
    I(6+2x+x^2+x^4)= 6x+x^2+\frac{x^3}{3}+\frac{x^5}{5}.
    \end{equation}
\end{example}

Remarquez que, étant donné que dans la définition de $I$ nous avons décidé d'intégrer entre zéro et $x$, tous les polynômes dans $\mathcal{P}_\eR^{k+1}$ qui sont l'image par $I$ d'un polynôme de $\mathcal{P}_\eR^{k}$ ont $a_0=0$. Cela veut dire que nous pouvons générer toute l'image de $I$ en utilisant un sous-ensemble de la base canonique de $\mathcal{P}_\eR^{k+1}$,  en particulier $\mathcal{B}_1=\{x\mapsto x^j \,|\, j=1, \ldots, k\}\subset \mathcal{B}$ nous suffira. Cela n'est guère surprenant, parce que l'image par une application linéaire d'un espace vectoriel de dimension finie ne peut pas être un espace de dimension supérieure.
\end{description}

Les applications de dérivation et intégration correspondent évidemment à des applications linéaires de $\mathcal{P}_\eR$ dans lui-même.

L'espace de tous les polynômes étant de dimension infinie, il peut servir de contre-exemple assez simple. Dans la sous-section~\ref{SubSecPOlynomesCE}, nous verrons que toutes les normes ne sont pas équivalentes sur l'espace des polynômes.

%+++++++++++++++++++++++++++++++++++++++++++++++++++++++++++++++++++++++++++++++++++++++++++++++++++++++++++++++++++++++++++ 
\section{Théorème de Sylvester}
%+++++++++++++++++++++++++++++++++++++++++++++++++++++++++++++++++++++++++++++++++++++++++++++++++++++++++++++++++++++++++++

\begin{definition}[Signature\cite{BIBooXOWGooAPWTfT}]       \label{DEFooWDCLooDkRYLK}
    Soit une forme quadratique\footnote{Définition \ref{DefBSIoouvuKR}.} \( Q\) sur un espace vectoriel \( E\) de dimension finie \( n\). L'\defe{indice d'inertie}{indice d'inertie} de \( q\) est le nombre
    \begin{equation}
        q=\max\{ \dim(F)\tq Q(v)<0\,\forall v\in F\setminus\{ 0 \} \}.
    \end{equation}
    Nous définissons aussi
    \begin{equation}
        p=\max\{ \dim(G)\tq Q(v)>0\,\forall v\in G\setminus\{ 0 \} \}.
    \end{equation}
    Le couple \( (p,q)\) est la \defe{signature}{signature!forme quadratique} de \( Q\).
\end{definition}

\begin{theorem}[de Sylvester\cite{BIBooXOWGooAPWTfT}]   \label{ThoQFVsBCk}
    Soit $Q$ une forme quadratique réelle de signature\footnote{Définition \ref{DEFooWDCLooDkRYLK}.} \( (p,q)\). Alors pour toute base orthonormée \( \{ e_i \}\) de \( \eR^{p+q}\) nous avons les points suivants.
    \begin{enumerate}
        \item
            Les nombres \( p\) et \( q\) sont donnée par 
    \begin{subequations}
        \begin{align}
            p&=\Card\{ i\tq Q(e_i)>0 \}\\
            q&=\Card\{ i\tq Q(e_i)<0 \}.
        \end{align}
    \end{subequations}
\item
    Le rang de \( Q\) est \( p+q\).
\item
    Si \( A\) est la matrice de \( Q\) dans une base, alors il existe une matrice inversible \( P\) telle que
    \begin{equation}
        P^tAP=\begin{pmatrix}
            -\mtu_q    &       &       \\
                &   \mtu_p    &       \\
                &       &   0
        \end{pmatrix}.
    \end{equation}
    \end{enumerate}
\end{theorem}
\index{théorème!Sylvester}
\index{rang}

\begin{proposition}[\cite{BIBooXOWGooAPWTfT}]
    Deux formes quadratiques sont équivalentes si et seulement si elles ont même signature.
\end{proposition}

\index{matrice!semblables}
\index{forme!quadratique}

%--------------------------------------------------------------------------------------------------------------------------- 
\subsection{Invariance de la trace}
%---------------------------------------------------------------------------------------------------------------------------

\begin{proposition}[\cite{MonCerveau}]      \label{PROPooRMYQooWkEpJJ}
    Soit une application linéaire \( f\). Si la matrice de \( f\) dans une base est \( A\) et est \( B\) dans une autre base, alors
    \begin{equation}
        \trace(A)=\trace(B).
    \end{equation}
\end{proposition}

\begin{proof}
    Les matrices \( A\) et \( B\) sont liées par la proposition \ref{PROPooNZBEooWyCXTw} : \( B=Q^{-1}AQ\) où \( Q\) est la matrice qui lie les vecteurs des deux bases. L'invariance cyclique de la trace donnée en le lemme \ref{LEMooUXDRooWZbMVN} implique que
    \begin{equation}
        \trace(B)=\trace(Q^{-1}AQ)=\trace(QQ^{-1}A)=\trace(A).
    \end{equation}
\end{proof}

%+++++++++++++++++++++++++++++++++++++++++++++++++++++++++++++++++++++++++++++++++++++++++++++++++++++++++++++++++++++++++++
\section{Dualité}
%+++++++++++++++++++++++++++++++++++++++++++++++++++++++++++++++++++++++++++++++++++++++++++++++++++++++++++++++++++++++++++


\begin{definition}  \label{DefJPGSHpn}
    Soit \( E\) un espace vectoriel sur \( \eK\).

    Une \defe{forme linéaire}{forme linéaire} sur \( E \) est une application linéaire de \( E \) sur son corps de base \( \eK\).

    Le \defe{dual algébrique}{dual algébrique} de \( E\), noté \( E^*\), l'ensemble des formes linéaires sur \( E\). Ainsi, \( E^* = \GL(E,\eK)\).
\end{definition}

Nous verrons plus tard qu'en dimension infinie, les applications linéaires ne sont pas toujours continues. Nous définirons donc aussi un concept de dual topologique. Voir la proposition~\ref{PROPooQZYVooYJVlBd}, la remarque~\ref{RemOAXNooSMTDuN} et la définition~\ref{DEFooKSDFooGIBtrG}.

\begin{definition}      \label{DEFooTMSEooZFtsqa}
    Si \( E\) est un espace vectoriel et si \( \{ e_i \}\) est une base de \( E\), alors nous définissons la \defe{base duale}{base!duale} de \( E^*\) par
    \begin{equation}
        e_i^*(e_j)=\delta_{ij}
    \end{equation}
    est sa prolongation par linéarité.
\end{definition}
Notons que si \( v\in E\) est un vecteur, ça n'a aucun sens à priori de parler de \( v^*\). Il s'agit bien de définir \emph{toute} la base \( \{ e_i^* \}\) à partir de toute la base \( \{ e_i \}\).

%---------------------------------------------------------------------------------------------------------------------------
\subsection{Orthogonal}
%---------------------------------------------------------------------------------------------------------------------------

\begin{definition}      \label{DEFooEQSMooHVzbfz}
    Soit \( E\), un espace vectoriel, et \( F\) une sous-espace de \( E\). L'\defe{orthogonal}{orthogonal!sous-espace} de \( F\) est la partie \( F^{\perp}\subset E^*\) donnée par
    \begin{equation}    \label{Eqiiyple}
        F^{\perp}=\{ \alpha\in E^*\tq \forall x\in F,\alpha(x)=0 \}.
    \end{equation}
\end{definition}

Cette définition d'orthogonal via le dual n'est pas du pur snobisme. En effet, la définition «usuelle» qui ne parle pas de dual,
\begin{equation}
    F^{\perp}=\{ y\in E\tq \forall x\in F,y\cdot x=0 \},
\end{equation}
demande la donnée d'un produit scalaire. Évidemment dans le cas de \( \eR^n\) munie du produit scalaire usuel et de l'identification usuelle entre \( \eR^n\) et \( (\eR^n)^*\) via une base, les deux notions d'orthogonal coïncident.

La définition~\ref{DEFooEQSMooHVzbfz}, au contraire, est intrinsèque : elle ne dépend que de la structure d'espace vectoriel.

Si \( B\subset E^*\), on note \( B^o\)\nomenclature[G]{\( B^o\)}{orthogonal dans le dual} son orthogonal :
\begin{equation}
    B^o=\{ x\in E\tq \omega(x)=0\,\forall \omega\in B \}.
\end{equation}
Notons qu'on le note \( B^o\) et non \( B^{\perp}\) parce qu'on veut un peu s'abstraire du fait que \( (E^*)^*=E\). Du coup on impose que \( B\) soit dans un dual et on prend une notation précise pour dire qu'on remonte au pré-dual et non qu'on va au dual du dual.

\begin{proposition} \label{PropXrTDIi}
    Soient un espace vectoriel \( E\) et un sous-espace vectoriel \( F\). Nous avons
    \begin{equation}
        \dim F+\dim F^{\perp}=\dim E.
    \end{equation}
\end{proposition}

\begin{proof}
    Soit \( \{ e_1,\ldots, e_p \}\) une base de \( F\) que nous complétons en une base \( \{ e_1,\ldots, e_n \}\) de \( E\) par le théorème~\ref{ThonmnWKs}. Soit \( \{ e_1^*,\ldots, e^*_n \}\) la base duale. Alors nous prouvons que \( \{ e^*_{p+1},\ldots, e_n^* \}\) est une base de \( F^{\perp}\).

    Déjà c'est une partie libre en tant que partie d'une base.

    Ensuite ce sont des éléments de \( F^{\perp}\) parce que si \( i\leq p\) et si \( k\geq 1\), nous avons \( e^*_{p+k}(e_i)=0\); donc oui, \( e^*_{p+k}\in F^{\perp}\).

    Enfin \( F^{\perp}\subset\Span\{ e_{k}^*, k \in \{p+1, \dots, n\}\}\) parce que si \( \omega=\sum_{k=1}^n\omega_ke_k^*\), alors \( \omega(e_i)=\omega_i\), mais nous savons que si \( \omega\in F^{\perp}\), alors \( \omega(e_i)=0\) pour \( i\leq p\). Donc \( \omega=\sum_{k=p+1}^n\omega_ke^*_k\).
\end{proof}

La proposition \ref{PROPooNITTooCYcrrT} donnera une version plus terre à terre de la proposition \ref{PropXrTDIi} en disant que si nous avons un produit scalaire, alors \( V=F\oplus F^{\perp}\) où \( F^{\perp}\) est cette fois défini comme l'orthogonal pour le produit scalaire.

%---------------------------------------------------------------------------------------------------------------------------
\subsection{Transposée : pas d'approche naïve}
%---------------------------------------------------------------------------------------------------------------------------
\label{SUBSECooGPXVooEYwIiJ}

Il est légitime, si \( t\colon E\to E\) est une application linéaire, de dire que sa transposée soit l'application linéaire \( t^t\colon E\to E\) dont la matrice est la matrice transposée de celle de \( t\). Lorsque nous travaillons sur \( \eR^n\) muni de la base canonique, cela ne pose pas de problèmes et nous pouvons écrire des égalités du type \( \langle x, Ay\rangle =\langle A^tx, y\rangle \).

\begin{proposition}[Matrice transposée et produit scalaire]     \label{PROPooNARVooEuhweD}
    Soit une matrice réelle \( A\). En utilisant l'application linéaire associée\footnote{Définition \ref{DEFooJVOAooUgGKme}. Ici nous considérons la base canonique sur \( \eR^n\)} \( f_A\colon \eR^n\to \eR^n\), nous avons
    \begin{equation}
        x\cdot f_A(y)=f_{A^t}(x)\cdot y.
    \end{equation}
    Cette formule est souvent écrire \( x\cdot Ay=A^tx\cdot y\) ou \( \langle x, Ay\rangle =\langle A^tx, y\rangle \).
\end{proposition}

\begin{proof}
    Il s'agit d'une calcul utilisant la formule \eqref{EQooBVGHooJhFbMs} et le produit scalaire \eqref{EQooFITHooEXDCGd} :
    \begin{equation}
        x\cdot f_A(y)=\sum_ix_i\big( f_A(y) \big)_i=\sum_ix_i\sum_jA_{ij}y_j=\sum_{ij}A^t_{ji}x_iy_j=\sum_jf_{A^t}(x)_jy_j=f_{A^t}(x)\cdot y.
    \end{equation}
\end{proof}

Hélas nous allons voir que cette façon de définir une transposée est mauvaise.

Soit une application linéaire \( t\colon E\to E\) de matrice \( A\) dans la base \( \{ e_i \}_{i=1,\ldots, n}\) et de matrice \( B\) dans la base \( \{ f_{\alpha} \}_{\alpha=1,\ldots, n}\).

Nous nommons \( t_1\) l'application linéaire associée à \( A^t\) dans la base \( \{ e_i \}\) et \( t_2\) l'application linéaire associée à la matrice \( B^t\) dans la base \( \{ f_{\alpha} \}\). Définir la transposée d'une application linéaire comme étant l'application linéaire associée à la transposée de sa matrice ne sera une bonne définition que si \( t_1=t_2\).

La première chose facile à voir est
\begin{equation}        \label{EQooAMHPooUQEkJo}
    t_1(e_i)_j=\sum_k(A^t)_{jk}(e_i)_k=A^t_{ji}=A_{ij}.
\end{equation}
Pour calculer \( t_2(e_i)_j\), c'est un peu plus laborieux :
\begin{subequations}
    \begin{align}
        t_2(e_i)&=\sum_{\alpha}Q_{\alpha i}^{-1} t_2(f_\alpha)=\sum_{\beta\gamma\alpha}Q_{\alpha i}^{-1}B^t_{\gamma\beta}\underbrace{(t_{\alpha})_{\beta}}_{\delta_{\alpha\beta}}f_{\gamma}=\sum_{\beta\gamma}Q_{\beta i}^{-1}B^t_{\gamma\beta}f_{\gamma}\\
        &=(B^tQ^{-1})_{\gamma i}Q_{j\gamma}e_j\\
        &=\sum_j(QB^tQ^{-1})_{ji}e_j.
    \end{align}
\end{subequations}
Donc \( t_2(e_i)_j=(QB^tQ^{-1})_{ji}\). En tenant compte du fait que \( B=Q^{-1}AQ\) nous avons
\begin{equation}
    t_2(e_i)_j=(QQ^tA^t(Q^{-1})^tQ^{-1})_{ji}.
\end{equation}
Cela est égal à l'expression \eqref{EQooAMHPooUQEkJo} lorsque \( Q^t=Q^{-1}\). Nous voyons que confondre transposée d'une application linéaire avec transposée de la matrice associée n'est valable que si nous sommes certain de ne considérer que des changements de base par des matrices orthogonales.

Cela est la situation typique dans laquelle nous nous trouvons lorsque nous considérons des applications linéaires sur \( \eR^n\) muni de la base canonique et que nous n'avons aucune intention de changer de base, et encore moins de chercher une base non orthonormale. Cette situation est clairement la situation la plus courante.

\begin{example}[\cite{ooLIOMooBuCPUS}]
    Soit la base canonique \( \{ e_1,e_2 \}\) de \( \eR^2\). Nous considérons l'application linéaire \( t\colon \eR^2\to \eR^2\) définie par
    \begin{subequations}
        \begin{align}
            t(e_1)&=e_1\\
            t(e_2)&=0.
        \end{align}
    \end{subequations}
    La matrice de \( t\) dans cette base est
    \begin{equation}
        A=\begin{pmatrix}
            1    &   0    \\
            0    &   0
        \end{pmatrix}.
    \end{equation}
    Elle est symétrique : elle vérifie \( A^t=A\). Si nous comptions sur la transposée de matrice pour définir la transposée de \( t\), nous aurions \( t^t=t\).

    Soit maintenant la base \( f_1=e_1\), \( f_2=e_1+e_2\). Nous avons \( t(f_1)=f_1\) et
    \begin{equation}
        t(f_2)=t(e_1)+t(e_2)=e_1=f_1.
    \end{equation}
    Donc la matrice de \( t\) dans cette base est
    \begin{equation}
        B=\begin{pmatrix}
            1    &   1    \\
            0    &   0
        \end{pmatrix}.
    \end{equation}
    Et là, nous avons \( B^t\neq B\). Donc en comptant sur cette base pour définir la transposée de \( t\) nous aurions \( t^t\neq t\).
\end{example}

\begin{normaltext}      \label{NooMZVRooExWVKJ}
    Autrement dit, la façon «usuelle» de voir la transposée d'une application linéaire ne fonctionne dans les livres pour enfant uniquement parce qu'on n'y considère toujours \( \eR^n\) muni de la base canonique ou de bases orthonormées.

    Notons que nous avons tout de même les notions d'opérateur adjoint et autoadjoint pour parler d'application orthogonale sans passer par la transposée, voir~\ref{DEFooYKCSooURQDoS}.
\end{normaltext}

%---------------------------------------------------------------------------------------------------------------------------
\subsection{Transposée : la bonne approche}
%---------------------------------------------------------------------------------------------------------------------------

\begin{definition}      \label{DefooZLPAooKTITdd}
    Si \( f\colon E\to F\) est une application linéaire entre deux espaces vectoriels, la \defe{transposée}{transposée} est l'application \( f^t\colon F^*\to E^*\) donnée par
    \begin{equation}
        f^t(\omega)(x)=\omega\big( f(x) \big).
    \end{equation}
    pour tout \( \omega\in F^*\) et \( x\in E\).
\end{definition}

\begin{lemma}
    Soit \( E\) muni de la base \( \{ e_i \}\) et \( F\) muni de la base \( \{ g_i \}\) et une application \( f\colon E\to F\). Si \( A\) est la matrice de \( f\) dans ces bases, alors \( A^t\) est la matrice de \( f^t\) dans les bases \( \{ e^*_i \}\) et \( \{ g^*_i \}\) de \( E^*\) et \( F^*\).
\end{lemma}

\begin{proof}
    Nous allons montrer que les formes \( f^t(g^*_i)\) et \( \sum_k(A^t)_{ik}g^*_k\) sont égales en les appliquant à un vecteur.

    Par définition de la matrice d'une application linéaire dans une base,
    \begin{equation}
        f^t(g_i^*)=\sum_j(f^t)_{ij}e^*_j,
    \end{equation}
    et
    \begin{equation}
        f(e_k)=\sum_lA_{kl}g_l.
    \end{equation}
    Du coup, si \( x=\sum_kx_ke_k\), nous avons
    \begin{equation}    \label{EqCzwftH}
        f^t(g_i^*)x=\sum_{kl}x_kg_i^*A_{kl}g_l=\sum_{kl}x_kA_{kl}\delta_{il}=\sum_k x_kA_{ki}=\sum_k(A^t)_{ik}x_k.
    \end{equation}
    D'autre part,
    \begin{equation}    \label{EqWlQlrR}
        \sum_k(A^t)_{ik}g_k^*x=\sum_{kl}(A^t)_{ik}g^*_kx_le_l=\sum_k(A^t)_{ik}x_k.
    \end{equation}
    Le fait que \eqref{EqCzwftH} et \eqref{EqWlQlrR} donnent le même résultat prouve le lemme.
\end{proof}
En corolaire, les rangs de \( f\) et de \( f^t\) sont égaux parce que le rang est donné par la plus grande matrice carrée de déterminant non nul. Nous prouvons cependant ce résultat de façon plus intrinsèque.

\begin{lemma}[\href{http://gilles.dubois10.free.fr/algebre_lineaire/dualite.html}{Gilles Dubois}]   \label{LemSEpTcW}
    Si \( f\colon E\to F\) est une application linéaire, alors
    \begin{equation}
        \rang(f)=\rang(f^t).
    \end{equation}
\end{lemma}

\begin{proof}
    Nous posons \( \dim\ker(f)=p\) et donc \( \rang(f)=n-p\). Soit \( \{ e_1,\ldots, e_p \}\) une base de \( \ker(f)\) que l'on complète en une base \( \{ e_1,\ldots, e_n \}\) de \( E\). Nous considérons maintenant les vecteurs
    \begin{equation}
        g_i=f(e_{p+i})
    \end{equation}
    pour \( i=1,\ldots, n-p\). C'est-à-dire que les \( g_i\) sont les images des vecteurs qui ne sont pas dans le noyau de \( f\). Prouvons qu'ils forment une famille libre. Si
    \begin{equation}
        \sum_{k=1}^{n-p}a_kf(e_{p+k})=0,
    \end{equation}
    alors \( f\big( \sum_ka_ke_{p+k} \big)=0\), ce qui signifierait que \( \sum_ka_ke_{p+k}\) se trouve dans le noyau de \( f\), ce qui est impossible par construction de la base \( \{ e_i \}_{i=1,\ldots, n}\). Étant donné que les vecteurs \( g_1,\ldots, g_{n-p}\) sont libres, nous les complétons en une base
    \begin{equation}
        \{ \underbrace{g_1,\ldots, g_{n-p}}_{\text{images}},\underbrace{g_{n-p+1},\ldots, g_r}_{\text{complétion}} \}
    \end{equation}
    de \( F\).

    Nous prouvons maintenant que \( \rang(f^t)\geq n-p\) en montrant que les formes \( \{ g_i^* \}_{i=1,\ldots, n-p}\) forment une partie libre (et donc l'espace image de \( f^t\) est au moins de dimension \( n-p\)). Pour cela nous prouvons que \( f^t(g_i^*)=e^*_{i+p}\). En effet
    \begin{equation}
        f^t(g^*_i)e_k=g_i^*(fe_k),
    \end{equation}
    Si \( k=1,\ldots, p\), alors \( fe_k=0\) et donc \( g_i^*(fe_k)=0\); si \( k=p+l\) alors
    \begin{equation}
        f^t(g_i^*)e_k=g_i^*(fe_{k+l})=g^*_i(g_l)=\delta_{i,l}=\delta_{i,k-p}=\delta_{k,i+p}.
    \end{equation}
    Donc \( f^t(g_i^*)=e^*_{i+p}\). Cela prouve que les formes \( f^t(g_i^*)\) sont libres et donc que
    \begin{equation}
        \rang(f^t)\geq n-p=\rang(f).
    \end{equation}
    En appliquant le même raisonnement à \( f^t\) au lieu de \( f\), nous trouvons
    \begin{equation}
        \rang\big( (f^t)^t \big)\geq \rang(f^t)
    \end{equation}
    et donc, vu que \( (f^t)^t=f\), nous obtenons \( \rang(f)=\rang(f^t)\).

\end{proof}

\begin{proposition}[\cite{DualMarcSAge}]        \label{PropWOPIooBHFDdP}
    Si \( f\) est une application linéaire entre les espaces vectoriels \( E\) et \( F\), alors nous avons
    \begin{equation}
        \Image(f^t)=\ker(f)^{\perp}.
    \end{equation}
\end{proposition}

\begin{proof}
    Soient donc l'application \( f\colon E\to F\) et sa transposée \( f^t\colon F^*\to E^*\). Nous commençons par prouver que \( \Image(f^{t})\subset(\ker f)^{\perp}\). Pour cela nous prenons \( \omega\in \Image(f^t)\), c'est-à-dire \( \omega=\alpha\circ f\) pour un certain élément \( \alpha\in F^*\). Si \( z\in\ker(f)\), alors \( \omega(z)=(\alpha\circ f)(z)=0\), c'est-à-dire que \( \omega\in (\ker f)^{\perp}\).

    Pour prouver qu'il y a égalité, nous n'allons pas démontrer l'inclusion inverse, mais plutôt prouver que les dimensions sont égales. Après, on sait que si \( A\subset B\) et si \( \dim A=\dim B\), alors \( A=B\). Nous avons
    \begin{subequations}
        \begin{align}
            \dim\big( \Image(f^t) \big)&=\rang(f^t)\\
            &=\rang(f)  &\text{lemme~\ref{LemSEpTcW}}\\
            &=\dim(E)-\dim\ker(f)   &\text{théorème~\ref{ThoGkkffA}}\\
            &=\dim\big( (\ker f)^{\perp} \big)  &\text{proposition~\ref{PropXrTDIi}}.
        \end{align}
    \end{subequations}
\end{proof}

\begin{lemma}[\cite{ooEPEFooQiPESf}]
    Soit \( \eK\) un corps, \( E\) et \( F\) deux \( \eK\)-espaces vectoriels de dimension finie et une application linéaire \( f\colon E\to F\). L'application \( f\) est injective si et seulement si sa transposée\footnote{Définition~\ref{DefooZLPAooKTITdd}.} \( f^t\) est surjective.
\end{lemma}

\begin{proof}
    Supposons que \( f\) soit injective. Alors par le lemme~\ref{LEMooDAACooElDsYb}, il existe \( g\colon F\to E\) tel que \( g\circ f=\id|_E\). Nous avons alors aussi \( (g\circ f)^t=\id|_{E^*}\), mais \( (g\circ f)^t=f^t\circ g^t\), donc \( f^t\) est surjective.

    Inversement, nous supposons que \( f^t\colon F^*\to E^*\) est surjective. Alors en nous souvenant que \( E\) et \( F\) sont de dimension finie et en faisant jouer les identifications \( (f^t)^t=f\) et \( (E^*)^*=E\) nous savons qu'il existe \( s\colon E^*\to F^*\) tel que \( f^t\circ s=\id|_{E^*}\). En passant à la transposée,
    \begin{equation}
        s^t\circ f=\id|_{E},
    \end{equation}
    qui implique que \( f\) est injective.
\end{proof}

%---------------------------------------------------------------------------------------------------------------------------
\subsection{Polynômes de Lagrange}
%---------------------------------------------------------------------------------------------------------------------------

Soit \( E=\eR_n[X]\) l'ensemble des polynômes à coefficients réels de degré au plus \( n\). Soient les \( n+1\) réels distincts \( a_0,\ldots, a_n\). Nous considérons les formes linéaires associées \( f_i\in E^*\),
\begin{equation}
    f_i(P)=P(a_i).
\end{equation}
\begin{lemma}
    Ces formes forment une base de \( E^*\).
\end{lemma}

\begin{proof}
    Nous prouvons que l'orthogonal est réduit au nul :
    \begin{equation}
        \Span\{ f_0,\ldots, f_n \}^{\perp}=\{ 0 \}
    \end{equation}
    pour que la proposition~\ref{PropXrTDIi} conclue. Si \( P\in\Span\{ f_i \}^{\perp}\), alors \( f_i(P)=0\) pour tout \( i\), ce qui fait que \( P(a_i)=0\) pour tout \( i=0,\ldots, n\). Un polynôme de degré au plus \( n\) qui s'annule en \( n+1\) points est automatiquement le polynôme nul.
\end{proof}

Les \defe{polynômes de Lagrange}{Lagrange!polynôme}\index{polynôme!Lagrange} sont les polynômes de la base (pré)duale de la base \( \{ f_i \}\).

\begin{proposition}
    Les polynômes de Lagrange sont donnés par
    \begin{equation}
        P_i=\prod_{k\neq i}\frac{ X-a_k }{ a_i-a_k }.
    \end{equation}
\end{proposition}

\begin{proof}
    Il suffit de vérifier que \( f_j(P_i)=\delta_{ij}\). Nous avons
    \begin{equation}
        f_j(P_i)=P_i(a_j)=\prod_{k\neq i}\frac{ a_j-a_k }{ a_i-a_k }.
    \end{equation}
    Si \( j\neq i\) alors un des termes est nul. Si au contraire \( i=j\), tous les termes valent \( 1\).
\end{proof}

%---------------------------------------------------------------------------------------------------------------------------
\subsection{Dual de \texorpdfstring{$ \eM(n,\eK)$}{M(n,K)}}
%---------------------------------------------------------------------------------------------------------------------------

\begin{proposition}[\cite{KXjFWKA}]     \label{PropHOjJpCa}
    Soit \( \eK\), un corps. Les formes linéaires sur \( \eM(n,\eK)\) sont les applications de la forme
    \begin{equation}
        \begin{aligned}
            f_A\colon \eM(n,\eK)&\to \eK \\
            M&\mapsto \tr(AM).
        \end{aligned}
    \end{equation}
\end{proposition}
\index{trace!dual de \( \eM(n,\eK)\)}
\index{dual!de \( \eM(n,\eK)\)}


\begin{proof}
    Nous considérons l'application
    \begin{equation}
        \begin{aligned}
            f\colon \eM(n,\eK)&\to \eM(n,\eK)^* \\
            A&\mapsto f_A
        \end{aligned}
    \end{equation}
    et nous voulons prouver que c'est une bijection. Étant donné que nous sommes en dimension finie, nous avons égalité des dimensions de \( \eM(n,\eK)\) et \( \left(\eM(n,\eK)\right)^*\), et il suffit de prouver que \( f\) est injective. Soit donc \( A\) telle que \( f_A=0\). Nous l'appliquons à la matrice \( (E_{ij})_{kl}=\delta_{ik}\delta_{jl}\) :
    \begin{equation}
            0=f_A(E_{ij})
            =\sum_{k}(AE_{ij})_{kk}
            =\sum_{kl}A_{kl}(E_{ij})_{lk}
            =\sum_{kl}A_{kl}\delta_{il}\delta_{jk}
            =A_{ij}.
    \end{equation}
    Donc \( A=0\).
\end{proof}

\begin{corollary}[\cite{KXjFWKA}]
    Soit \( \eK\) un corps et \( \phi\in\eM(n,\eK)^*\) telle que pour tout \( M,N\in \eM(n,\eK)\) on ait
    \begin{equation}
        \phi(MN)=\phi(NM).
    \end{equation}
    Alors il existe \( \lambda\in \eK\) tel que \( \phi=\lambda\Tr\).
\end{corollary}
\index{trace!unicité pour la propriété de trace}

\begin{proof}
    La proposition~\ref{PropHOjJpCa} nous donne une matrice \( A\in \eM(n,\eK)\) telle que \( \phi=f_A\). L'hypothèse nous dit que \( f_A(MN)=f_A(NM)\), c'est-à-dire
    \begin{equation}
        \Tr(AMN)=\Tr(ANM)
    \end{equation}
    pour toutes matrices \( M, N\in \eM(n,\eK)\). L'invariance cyclique de la trace\footnote{Lemme~\ref{LEMooUXDRooWZbMVN}.} appliqué au membre de droite nous donne \( \Tr(AMN)=\Tr(MAN)\), ce qui signifie que
    \begin{equation}
        \Tr\big( (AM-MA)N \big)=0
    \end{equation}
    ou encore que \( f_{AM-MA}=0\), et ce, pour toute matrice \( M\). La fonction \( f\) étant injective nous en déduisons que la matrice \( A\) doit satisfaire
    \begin{equation}
        AM=MA
    \end{equation}
    pour tout \( M\in\eM(n,\eK)\). En particulier, en prenant pour \( M \) les fameuses matrices \( E_{ij}\) et en calculant un peu,
    \begin{equation}
        A_{li}\delta_{jm}=\delta_{il}A_{jm}
    \end{equation}
    pour tout \( i,j,l,m\). Cela implique que \( A_{ll}=A_{mm}\) pour tout \( l\) et \( m\) et que \( A_{jm}=0\) dès que \( j\neq m\). Il existe donc \( \lambda\in \eK\) tel que \( A=\lambda\mtu\). En fin de compte,
    \begin{equation}
        \phi(X)=f_{\lambda\mtu}(X)=\lambda\Tr(X).
    \end{equation}
\end{proof}

\begin{corollary}[\cite{KXjFWKA}]       \label{CorICUOooPsZQrg}
    Soit \( \eK\) un corps. Tout hyperplan de \( \eM(n,\eK)\) coupe \( \GL(n,\eK)\).
\end{corollary}
\index{groupe!linéaire!hyperplan}

\begin{proof}
    Soit \( \mH\) un hyperplan de \( \eM\). Il existe une forme linéaire \( \phi\) sur \( \eM(n,\eK)\) telle que \( \mH=\ker(\phi)\). Encore une fois la proposition~\ref{PropHOjJpCa} nous donne \( A\in \eM\) telle que \( \phi=f_A\); nous notons \( r\) le rang de \( A\). Par le lemme~\ref{LemZMxxnfM} nous avons \( A=PJ_rQ\) avec \( P,Q\in \GL(n,\eK)\) et
    \begin{equation}
        J_r=\begin{pmatrix}
            \mtu_r    &   0    \\
            0    &   0
        \end{pmatrix}.
    \end{equation}
    Pour tout \( M\in \eM\) nous avons
    \begin{equation}
        \phi(M)=\Tr(AM)=\Tr(PJ_rQM)=\Tr(J_rQMP),
    \end{equation}
la dernière égalité découlant de l'invariance cyclique de la trace\footnote{Lemme~\ref{LEMooUXDRooWZbMVN}.}. Ce que nous cherchons est \( M\in \GL(n,\eK)\) telle que \( \phi(M)=0\). Nous commençons par trouver \( N\in\GL(n,\eK)\) telle que \( \Tr(J_rN)=0\). Celle-là est facile : c'est
    \begin{equation}
        N=\begin{pmatrix}
            0    &   1    \\
            \mtu_{n-1}    &   0
        \end{pmatrix}.
    \end{equation}
    Les éléments diagonaux de \( J_rN\) sont tous nuls. Par conséquent en posant \( M=Q^{-1}NP^{-1}\) nous avons notre matrice inversible dans le noyau de \( \phi\).
\end{proof}
\index{hyperplan!de \( \eM(n,\eK)\)}

%+++++++++++++++++++++++++++++++++++++++++++++++++++++++++++++++++++++++++++++++++++++++++++++++++++++++++++++++++++++++++++ 
\section{Représentation de groupe}
%+++++++++++++++++++++++++++++++++++++++++++++++++++++++++++++++++++++++++++++++++++++++++++++++++++++++++++++++++++++++++++

\begin{definition}[Représentation]      \label{DEFooXVMSooXDIfZV}
    Soit un groupe \( G\) et un espace vectoriel \( V\). Nous disons qu'une application \( \rho\colon G\to \GL(V)\) est une \defe{représentation}{représentation} de \( G\) sur \( V\) si pour tout \( g,h\in G\) nous avons
    \begin{equation}
        \rho(g)\circ\rho(g)=\rho(gh).
    \end{equation}
    Très souvent, nous disons que la représentation est le couple \( (V,\rho)\).
\end{definition}

\begin{definition}
    Une représentation\footnote{Définition \ref{DEFooXVMSooXDIfZV}.} est \defe{fidèle}{représentation!fidèle} si elle est injective en tant que application \( G\to \GL(V)\). Ce ne sont pas chacun des \( \rho(g)\) qui doivent être injectifs. La dimension de \( V\) est le \defe{degré}{degré!d'une représentation} de la représentation \( (V,\rho)\).
\end{definition}

\begin{proposition}     \label{PROPooHNQOooSzeEFG}
    Soit un corps \( \eK\). Si \( G\) est un groupe dans \( \eM(n,\eK)\) (c'est-à-dire un groupe de matrices à coefficients dans \( \eK\)), alors l'application
    \begin{equation}
        \begin{aligned}
                \rho\colon G&\to \GL(\eK^n) \\
            A&\mapsto f_A 
        \end{aligned}
    \end{equation}
    où \( f_A\) est l'application linéaire associée à \( A\) est une représentation de \( G\).
\end{proposition}


